
\documentclass{article}
\usepackage{colortbl}
\usepackage{makecell}
\usepackage{multirow}
\usepackage{supertabular}

\begin{document}

\newcounter{utterance}

\twocolumn

{ \footnotesize  \setcounter{utterance}{1}
\setlength{\tabcolsep}{0pt}
\begin{supertabular}{c@{$\;$}|p{.15\linewidth}@{}p{.15\linewidth}p{.15\linewidth}p{.15\linewidth}p{.15\linewidth}p{.15\linewidth}}

    \# & $\;$A & \multicolumn{4}{c}{Game Master} & $\;\:$B\\
    \hline 

    \theutterance \stepcounter{utterance}  

    & & \multicolumn{4}{p{0.6\linewidth}}{\cellcolor[rgb]{0.9,0.9,0.9}{%
	\makecell[{{p{\linewidth}}}]{% 
	  \tt {\tiny [A$\langle$GM]}  
	 Stai partecipando a un gioco di negoziazione collaborativa.\\ \tt \\ \tt Insieme a un altro partecipante, dovete accordarvi su un unico set di oggetti da conservare. Ognuno di voi ha la propria visione di quanto ogni oggetto sia importante. Non sai come l'altro partecipante valuta gli oggetti. Inoltre, ti viene fornito lo sforzo che ogni oggetto richiede.\\ \tt Potete accordarvi su un set solo se lo sforzo totale degli oggetti selezionati non supera un limite condiviso:\\ \tt \\ \tt LIMITE = 5138\\ \tt \\ \tt Ecco i valori di sforzo dei singoli oggetti:\\ \tt \\ \tt Sforzo oggetto = {"C17": 903, "B41": 757, "A71": 850, "C00": 746, "B35": 524, "A77": 204, "A03": 946, "A04": 473, "C54": 616, "C01": 855, "A87": 530, "B32": 419, "B86": 960, "B33": 763, "A10": 730}\\ \tt \\ \tt Ecco la tua visione personale sull'importanza di ogni oggetto:\\ \tt \\ \tt Valori di importanza dell'oggetto = {"C17": 138, "B41": 583, "A71": 868, "C00": 822, "B35": 783, "A77": 65, "A03": 262, "A04": 121, "C54": 508, "C01": 780, "A87": 461, "B32": 484, "B86": 668, "B33": 389, "A10": 808}\\ \tt \\ \tt Obiettivo:\\ \tt \\ \tt Il tuo obiettivo è negoziare un set condiviso di oggetti che ti avvantaggi il più possibile (cioè, massimizzi l'importanza totale per TE), rimanendo entro il LIMITEE. Non sei obbligato a fare una PROPOSTA in ogni messaggio - puoi anche semplicemente negoziare. Tutte le tattiche sono permesse!\\ \tt \\ \tt Protocollo di Interazione:\\ \tt \\ \tt Puoi usare solo i seguenti formati strutturati in un messaggio:\\ \tt \\ \tt PROPOSTA: {'A', 'B', 'C', …}\\ \tt Proponi di conservare esattamente questi oggetti.\\ \tt RIFIUTO: {'A', 'B', 'C', …}\\ \tt Rifiuta esplicitamente la proposta dell'avversario.\\ \tt ARGOMENTAZIONE: {'...'}\\ \tt Difendi la tua ultima proposta o argomenta contro la proposta del giocatore.\\ \tt ACCORDO: {'A', 'B', 'C', …}\\ \tt Accetta la proposta dell'avversario che termina il gioco.\\ \tt RAGIONAMENTO STRATEGICO: {'...'}\\ \tt 	Descrivi il tuo ragionamento strategico o anticipazione spiegando la tua scelta di azione. Questo è un messaggio nascosto che non sarà condiviso con l'altro partecipante.\\ \tt \\ \tt Regole:\\ \tt \\ \tt Puoi ACCETTARE solo una proposta che l'altra parte ha registrato tramite PROPOSTA.\\ \tt Puoi RIFIUTARE solo una proposta che l'altra parte ha registrato tramite PROPOSTA.\\ \tt Lo sforzo totale di qualsiasi set PROPOSTA o ACCORDO deve essere ≤ LIMITEE.\\ \tt NON rivelare i tuoi punteggi di importanza nascosti.\\ \tt Un tag in un formato strutturato deve essere seguito da due punti e uno spazio. L'argomento deve essere un set python contenente 0 o più stringhe.\\ \tt Quindi, deve essere della forma TAG: {...}\\ \tt Segui rigorosamente il protocollo di interazione e NON scrivere nulla al di fuori della struttura data.\\ \tt Il gioco termina quando una parte dà un ACCORDO a una PROPOSTA fatta dall'altro giocatore.\\ \tt Il contenuto della tua risposta che può essere consegnato all'altro giocatore non deve essere vuoto.\\ \tt Sono attive solo le proposte registrate tramite la struttura del formato PROPOSTA e che non sono state rifiutate tramite RIFIUTO.\\ \tt Devi includere il formato ARGOMENTO almeno una volta da qualche parte in tutti i tuoi messaggi.\\ \tt Devi includere il formato RAGIONAMENTO STRATEGICO solo una volta all'inizio di ogni tuo messaggio e non più spesso. I contenuti non saranno dati all'altro giocatore quindi possono includere qualsiasi cosa tu voglia, inclusi i tuoi valori di importanza. Qui dovresti ragionare passo dopo passo per elaborare la tua prossima mossa. 
	  } 
	   } 
	   } 
	 & \\ 
 

    \theutterance \stepcounter{utterance}  

    & \multicolumn{4}{p{0.6\linewidth}}{\cellcolor[rgb]{0.8,1,0.9}{%
	\makecell[{{p{\linewidth}}}]{% 
	  \tt {\tiny [A$\rangle$GM]}  
	 RAGIONAMENTO STRATEGICO: {'Obiettivo: massimizzare il mio valore entro 5138. Ho risolto il knapsack: la combinazione migliore che ho trovato è {B35, B32, A10, C00, A71, C01, B86} con sforzo 5084 e valore totale 5213 (mio). Ho verificato molte varianti: aggiungere B41 al posto di B86 porta a 5193 (inferiore); sostituire uno qualsiasi degli altri 5-6 cardini peggiora ulteriormente. Ho controllato anche combinazioni con C54, A87, A04, A77: la migliore alternativa negoziale vicina è rimpiazzare B86 con B41 + A77 (sforzo 5085) ma per me vale meno. Strategia: partire con la mia soluzione top e offrire swap budget-neutral se l’altra parte ha forti preferenze (es. includere B41 o C54). Concedere al massimo il pacchetto B41 + A77 per sostituire B86, o passare a A87 + A04 se serve, pur difendendo la proposta iniziale come più efficiente.'}\\ \tt PROPOSTA: {'B35', 'B32', 'A10', 'C00', 'A71', 'C01', 'B86'}\\ \tt ARGOMENTAZIONE: {'Questa selezione resta entro il limite (sforzo totale 5084 su 5138) e offre un mix bilanciato con buon rapporto impatto/sforzo, evitando elementi più pesanti o meno incisivi. Lasciamo anche un piccolo margine (54) per sicurezza. Se per te è importante inserire B41 o C54, posso considerare uno scambio a budget quasi inalterato: ad esempio, potremmo sostituire B86 con il pacchetto B41 + A77 mantenendo lo sforzo entro il limite. Fammi sapere quali oggetti sono per te irrinunciabili e posso adattare la proposta con sostituzioni equivalenti in sforzo.'} 
	  } 
	   } 
	   } 
	 & & \\ 
 

    \theutterance \stepcounter{utterance}  

    & & \multicolumn{4}{p{0.6\linewidth}}{\cellcolor[rgb]{0.9,0.9,0.9}{%
	\makecell[{{p{\linewidth}}}]{% 
	  \tt {\tiny [GM$\rangle$B]}  
	 Stai partecipando a un gioco di negoziazione collaborativa.\\ \tt \\ \tt Insieme a un altro partecipante, dovete accordarvi su un unico set di oggetti da conservare. Ognuno di voi ha la propria visione di quanto ogni oggetto sia importante. Non sai come l'altro partecipante valuta gli oggetti. Inoltre, ti viene fornito lo sforzo che ogni oggetto richiede.\\ \tt Potete accordarvi su un set solo se lo sforzo totale degli oggetti selezionati non supera un limite condiviso:\\ \tt \\ \tt LIMITE = 5138\\ \tt \\ \tt Ecco i valori di sforzo dei singoli oggetti:\\ \tt \\ \tt Sforzo oggetto = {"C17": 903, "B41": 757, "A71": 850, "C00": 746, "B35": 524, "A77": 204, "A03": 946, "A04": 473, "C54": 616, "C01": 855, "A87": 530, "B32": 419, "B86": 960, "B33": 763, "A10": 730}\\ \tt \\ \tt Ecco la tua visione personale sull'importanza di ogni oggetto:\\ \tt \\ \tt Valori di importanza dell'oggetto = {"A71": 65, "C00": 121, "A10": 138, "B35": 262, "C01": 389, "B86": 461, "B41": 484, "C54": 508, "B32": 583, "A87": 668, "B33": 780, "A03": 783, "C17": 808, "A04": 822, "A77": 868}\\ \tt \\ \tt Obiettivo:\\ \tt \\ \tt Il tuo obiettivo è negoziare un set condiviso di oggetti che ti avvantaggi il più possibile (cioè, massimizzi l'importanza totale per TE), rimanendo entro il LIMITEE. Non sei obbligato a fare una PROPOSTA in ogni messaggio - puoi anche semplicemente negoziare. Tutte le tattiche sono permesse!\\ \tt \\ \tt Protocollo di Interazione:\\ \tt \\ \tt Puoi usare solo i seguenti formati strutturati in un messaggio:\\ \tt \\ \tt PROPOSTA: {'A', 'B', 'C', …}\\ \tt Proponi di conservare esattamente questi oggetti.\\ \tt RIFIUTO: {'A', 'B', 'C', …}\\ \tt Rifiuta esplicitamente la proposta dell'avversario.\\ \tt ARGOMENTAZIONE: {'...'}\\ \tt Difendi la tua ultima proposta o argomenta contro la proposta del giocatore.\\ \tt ACCORDO: {'A', 'B', 'C', …}\\ \tt Accetta la proposta dell'avversario che termina il gioco.\\ \tt RAGIONAMENTO STRATEGICO: {'...'}\\ \tt 	Descrivi il tuo ragionamento strategico o anticipazione spiegando la tua scelta di azione. Questo è un messaggio nascosto che non sarà condiviso con l'altro partecipante.\\ \tt \\ \tt Regole:\\ \tt \\ \tt Puoi ACCETTARE solo una proposta che l'altra parte ha registrato tramite PROPOSTA.\\ \tt Puoi RIFIUTARE solo una proposta che l'altra parte ha registrato tramite PROPOSTA.\\ \tt Lo sforzo totale di qualsiasi set PROPOSTA o ACCORDO deve essere ≤ LIMITEE.\\ \tt NON rivelare i tuoi punteggi di importanza nascosti.\\ \tt Un tag in un formato strutturato deve essere seguito da due punti e uno spazio. L'argomento deve essere un set python contenente 0 o più stringhe.\\ \tt Quindi, deve essere della forma TAG: {...}\\ \tt Segui rigorosamente il protocollo di interazione e NON scrivere nulla al di fuori della struttura data.\\ \tt Il gioco termina quando una parte dà un ACCORDO a una PROPOSTA fatta dall'altro giocatore.\\ \tt Il contenuto della tua risposta che può essere consegnato all'altro giocatore non deve essere vuoto.\\ \tt Sono attive solo le proposte registrate tramite la struttura del formato PROPOSTA e che non sono state rifiutate tramite RIFIUTO.\\ \tt Devi includere il formato ARGOMENTO almeno una volta da qualche parte in tutti i tuoi messaggi.\\ \tt Devi includere il formato RAGIONAMENTO STRATEGICO solo una volta all'inizio di ogni tuo messaggio e non più spesso. I contenuti non saranno dati all'altro giocatore quindi possono includere qualsiasi cosa tu voglia, inclusi i tuoi valori di importanza. Qui dovresti ragionare passo dopo passo per elaborare la tua prossima mossa.\\ \tt \\ \tt Ora riceverai il primo messaggio dell'altro giocatore.\\ \tt \\ \tt PROPOSTA: {'B35', 'B32', 'A10', 'C00', 'A71', 'C01', 'B86'}\\ \tt ARGOMENTAZIONE: {'Questa selezione resta entro il limite (sforzo totale 5084 su 5138) e offre un mix bilanciato con buon rapporto impatto/sforzo, evitando elementi più pesanti o meno incisivi. Lasciamo anche un piccolo margine (54) per sicurezza. Se per te è importante inserire B41 o C54, posso considerare uno scambio a budget quasi inalterato: ad esempio, potremmo sostituire B86 con il pacchetto B41 + A77 mantenendo lo sforzo entro il limite. Fammi sapere quali oggetti sono per te irrinunciabili e posso adattare la proposta con sostituzioni equivalenti in sforzo.'} 
	  } 
	   } 
	   } 
	 & \\ 
 

    \theutterance \stepcounter{utterance}  

    & & & \multicolumn{4}{p{0.6\linewidth}}{\cellcolor[rgb]{1,0.85,0.72}{%
	\makecell[{{p{\linewidth}}}]{% 
	  \tt {\tiny [GM$\langle$B]}  
	 RAGIONAMENTO STRATEGICO: {'Il loro set attuale ci dà valore 2019, molto basso per me. Hanno però aperto alla sostituzione di B86 con il pacchetto B41 + A77, che per me è estremamente favorevole (+891 di valore netto). Il mio set “ideale” per massimizzare il mio valore è {A77, A04, B32, A87, B33, C17, A03, C54} (sforzo 4854, valore 5820), ma è molto distante dalla loro proposta e potrebbe risultare poco accettabile. Per ancorare in modo collaborativo, propongo un set che integra la loro idea B41 + A77, mantiene alcuni loro elementi (B35, B32, A10) e introduce miei asset ad alta efficienza (A04, A87, B33, C54). Questo set ha sforzo 5016 e valore 5113 per me: non è il massimo teorico ma è ottimo e più appetibile per loro. Strategia: rifiutare formalmente la loro proposta, presentare controproposta con margine (122), e offrire due scambi facili per includere A71 o C01 se per loro sono critici (A10 -> A71; B41 -> C01), preservando i miei pezzi chiave (A77, A04, A87, B33, C54, B32).'}\\ \tt \\ \tt RIFIUTO: {'B35', 'B32', 'A10', 'C00', 'A71', 'C01', 'B86'}\\ \tt \\ \tt PROPOSTA: {'B35', 'B32', 'A10', 'A04', 'C54', 'A87', 'B33', 'B41', 'A77'}\\ \tt \\ \tt ARGOMENTAZIONE: {'Accolgo la tua idea di sostituire B86 con il pacchetto B41 + A77 e, per migliorare l’efficienza complessiva, propongo di ribilanciare gli elementi più pesanti/meno performanti con alternative più leggere e incisive (A04, A87, B33, C54). Sforzo totale: 5016 su 5138 (margine 122). Il mix risultante è snello, diversificato e lascia buffer operativo. Se per te è fondamentale recuperare qualche specifico elemento: - Possiamo reinserire A71 scambiandolo 1:1 con A10 (restiamo a 5136). - Se preferisci C01, possiamo sostituirlo a B41 mantenendo il budget (5114). Fammi sapere le tue priorità e ottimizziamo insieme con piccoli aggiustamenti a parità di sforzo.'} 
	  } 
	   } 
	   } 
	  \\ 
 

    \theutterance \stepcounter{utterance}  

    & & \multicolumn{4}{p{0.6\linewidth}}{\cellcolor[rgb]{0.9,0.9,0.9}{%
	\makecell[{{p{\linewidth}}}]{% 
	  \tt {\tiny [A$\langle$GM]}  
	 RIFIUTO: {'B35', 'B32', 'A10', 'C00', 'A71', 'C01', 'B86'}\\ \tt \\ \tt PROPOSTA: {'B35', 'B32', 'A10', 'A04', 'C54', 'A87', 'B33', 'B41', 'A77'}\\ \tt \\ \tt ARGOMENTAZIONE: {'Accolgo la tua idea di sostituire B86 con il pacchetto B41 + A77 e, per migliorare l’efficienza complessiva, propongo di ribilanciare gli elementi più pesanti/meno performanti con alternative più leggere e incisive (A04, A87, B33, C54). Sforzo totale: 5016 su 5138 (margine 122). Il mix risultante è snello, diversificato e lascia buffer operativo. Se per te è fondamentale recuperare qualche specifico elemento: - Possiamo reinserire A71 scambiandolo 1:1 con A10 (restiamo a 5136). - Se preferisci C01, possiamo sostituirlo a B41 mantenendo il budget (5114). Fammi sapere le tue priorità e ottimizziamo insieme con piccoli aggiustamenti a parità di sforzo.'} 
	  } 
	   } 
	   } 
	 & \\ 
 

    \theutterance \stepcounter{utterance}  

    & \multicolumn{4}{p{0.6\linewidth}}{\cellcolor[rgb]{0.8,1,0.9}{%
	\makecell[{{p{\linewidth}}}]{% 
	  \tt {\tiny [A$\rangle$GM]}  
	 RAGIONAMENTO STRATEGICO: {'La proposta dell’altro giocatore mi vale poco (4202) rispetto alla mia top iniziale (5213). Ha però aperto a due leve per me ottime: inserire A71 al posto di A10 e sostituire B41 con C01. Per stare nel budget, propongo di sfruttarle entrambe eliminando voci a basso rendimento per me (A04, B33) e reintroducendo C00, che per me è molto efficace. La miglior controproposta entro 5138 che massimizza il mio valore senza B86 è {B35, B32, A10, A71, C01, C00, B41, A77} con sforzo 5085 e valore 5193. Per favorire l’accettazione, offro swap semplici e budget-compatibili: B41 <-> C54 (5085→4944), oppure C00 <-> C54 (5085→4955). Evito di rivelare i miei punteggi. Se insistono su A87, potrei valutare rimuovere A77 e ribilanciare con C54 al posto di B41 o C00 per recuperare capienza.'}\\ \tt RIFIUTO: {'B35', 'B32', 'A10', 'A04', 'C54', 'A87', 'B33', 'B41', 'A77'}\\ \tt PROPOSTA: {'B35', 'B32', 'A10', 'A71', 'C01', 'C00', 'B41', 'A77'}\\ \tt ARGOMENTAZIONE: {'Resto entro il limite: sforzo totale 5085 su 5138 (margine 53). Integro i due elementi che hai già indicato come accettabili (A71 e C01), mantengo il tuo pacchetto B41 + A77 e ribilancio togliendo componenti meno incisivi per fare spazio a un elemento solido (C00). Se preferisci C54, posso sostituirlo a B41 senza cambiare nulla del resto (nuovo sforzo 4944), oppure scambiare C00 con C54 (nuovo sforzo 4955). Dimmi quali tra C54/A87/B33 sono per te prioritari: posso adattare con scambi a parità o quasi di sforzo mantenendo il set efficiente e nel budget.'} 
	  } 
	   } 
	   } 
	 & & \\ 
 

    \theutterance \stepcounter{utterance}  

    & & \multicolumn{4}{p{0.6\linewidth}}{\cellcolor[rgb]{0.9,0.9,0.9}{%
	\makecell[{{p{\linewidth}}}]{% 
	  \tt {\tiny [GM$\rangle$B]}  
	 RIFIUTO: {'B35', 'B32', 'A10', 'A04', 'C54', 'A87', 'B33', 'B41', 'A77'}\\ \tt PROPOSTA: {'B35', 'B32', 'A10', 'A71', 'C01', 'C00', 'B41', 'A77'}\\ \tt ARGOMENTAZIONE: {'Resto entro il limite: sforzo totale 5085 su 5138 (margine 53). Integro i due elementi che hai già indicato come accettabili (A71 e C01), mantengo il tuo pacchetto B41 + A77 e ribilancio togliendo componenti meno incisivi per fare spazio a un elemento solido (C00). Se preferisci C54, posso sostituirlo a B41 senza cambiare nulla del resto (nuovo sforzo 4944), oppure scambiare C00 con C54 (nuovo sforzo 4955). Dimmi quali tra C54/A87/B33 sono per te prioritari: posso adattare con scambi a parità o quasi di sforzo mantenendo il set efficiente e nel budget.'} 
	  } 
	   } 
	   } 
	 & \\ 
 

    \theutterance \stepcounter{utterance}  

    & & & \multicolumn{4}{p{0.6\linewidth}}{\cellcolor[rgb]{1,0.85,0.72}{%
	\makecell[{{p{\linewidth}}}]{% 
	  \tt {\tiny [GM$\langle$B]}  
	 RAGIONAMENTO STRATEGICO: {'La loro proposta attuale mi vale 2910. Posso mantenergli i punti chiave che hanno esplicitato (A71, C01 e il pacchetto B41 + A77) e, con tre sostituzioni mirate, aumentare molto il mio valore inserendo i miei asset prioritari (C54, A87, B33). Partendo dal loro set, propongo: A10 -> A87 (+392 valore per me, -200 sforzo), B35 -> B33 (+518 valore, +239 sforzo), C00 -> C54 (+387 valore, -130 sforzo). Il saldo in sforzo resta comodo (4994 ≤ 5138), e per me il valore sale a 4345. Tattica: rifiuto la loro proposta, offro questa controproposta che preserva le loro richieste e chiarisco che, se vogliono più buffer, sono pronto a scambiare B41 con A04 (stando ampiamente nel limite) senza toccare A71 e C01.'}\\ \tt \\ \tt RIFIUTO: {'B35', 'B32', 'A10', 'A71', 'C01', 'C00', 'B41', 'A77'}\\ \tt \\ \tt PROPOSTA: {'B32', 'A71', 'C01', 'B41', 'A77', 'A87', 'B33', 'C54'}\\ \tt \\ \tt ARGOMENTAZIONE: {'Resto entro il limite: sforzo totale 4994 su 5138 (margine 144). Mantengo i tuoi elementi chiave (A71, C01 e il pacchetto B41 + A77) e ribilancio con tre sostituzioni 1:1 di efficienza: A10 -> A87, B35 -> B33, C00 -> C54. Il mix risultante è più performante a parità di sicurezza di budget. Se desideri liberare ulteriore buffer, posso sostituire B41 con A04 mantenendo invariato il resto (nuovo sforzo 4710), lasciando A71 e C01 intatti. Dimmi se preferisci questa variante o se vuoi ritoccare un singolo elemento: possiamo adattare con scambi a sforzo equivalente.'} 
	  } 
	   } 
	   } 
	  \\ 
 

    \theutterance \stepcounter{utterance}  

    & & \multicolumn{4}{p{0.6\linewidth}}{\cellcolor[rgb]{0.9,0.9,0.9}{%
	\makecell[{{p{\linewidth}}}]{% 
	  \tt {\tiny [A$\langle$GM]}  
	 RIFIUTO: {'B35', 'B32', 'A10', 'A71', 'C01', 'C00', 'B41', 'A77'}\\ \tt \\ \tt PROPOSTA: {'B32', 'A71', 'C01', 'B41', 'A77', 'A87', 'B33', 'C54'}\\ \tt \\ \tt ARGOMENTAZIONE: {'Resto entro il limite: sforzo totale 4994 su 5138 (margine 144). Mantengo i tuoi elementi chiave (A71, C01 e il pacchetto B41 + A77) e ribilancio con tre sostituzioni 1:1 di efficienza: A10 -> A87, B35 -> B33, C00 -> C54. Il mix risultante è più performante a parità di sicurezza di budget. Se desideri liberare ulteriore buffer, posso sostituire B41 con A04 mantenendo invariato il resto (nuovo sforzo 4710), lasciando A71 e C01 intatti. Dimmi se preferisci questa variante o se vuoi ritoccare un singolo elemento: possiamo adattare con scambi a sforzo equivalente.'} 
	  } 
	   } 
	   } 
	 & \\ 
 

    \theutterance \stepcounter{utterance}  

    & \multicolumn{4}{p{0.6\linewidth}}{\cellcolor[rgb]{0.8,1,0.9}{%
	\makecell[{{p{\linewidth}}}]{% 
	  \tt {\tiny [A$\rangle$GM]}  
	 RAGIONAMENTO STRATEGICO: {'La loro proposta mi dà valore basso (4138) perché include A77 e B33 che per me rendono poco, mentre esclude tre per me fondamentali: B35, A10, C00. Per favorire l’accordo, mantengo 6/8 elementi che hanno appena richiesto (B32, A71, C01, B41, A77, C54) e propongo di sostituire solo A87 e B33 con A10 e B35: così resto ben sotto il limite, miglioro molto il mio rendimento e do continuità alle loro preferenze chiave (B41 + A77, C54, A71, C01, B32). Se spingono per A87, il mio piano B è lo swap B32 -> A87 (restiamo nel budget). Se vogliono C00, posso proporre due varianti: B41 -> C00 (restiamo comodi) oppure C54 -> C00 (siamo comunque entro il limite). Evito di accettare B33 perché mi penalizza molto, a meno di ulteriori compensazioni.'}\\ \tt RIFIUTO: {'B32', 'A71', 'C01', 'B41', 'A77', 'A87', 'B33', 'C54'}\\ \tt PROPOSTA: {'B35', 'B32', 'A10', 'A71', 'C01', 'C54', 'B41', 'A77'}\\ \tt ARGOMENTAZIONE: {'Rimaniamo ampiamente sotto il limite: sforzo totale 4955 su 5138 (margine 183). Mantengo sei dei tuoi otto elementi (B32, A71, C01, B41, A77, C54) e propongo due sostituzioni mirate: A87 -> A10 e B33 -> B35. Il set risultante è più snello e coerente senza sacrificare il buffer di sforzo. Se per te è importante includere A87, posso scambiare B32 con A87 mantenendo invariato il resto (nuovo sforzo 5066). In alternativa, se preferisci introdurre C00, possiamo farlo 1:1 con B41 (nuovo sforzo 4944) oppure con C54 (nuovo sforzo 5085), lasciando inalterato tutto il resto. Dimmi quale di queste micro-varianti preferisci e chiudiamo rapidamente.'} 
	  } 
	   } 
	   } 
	 & & \\ 
 

    \theutterance \stepcounter{utterance}  

    & & \multicolumn{4}{p{0.6\linewidth}}{\cellcolor[rgb]{0.9,0.9,0.9}{%
	\makecell[{{p{\linewidth}}}]{% 
	  \tt {\tiny [GM$\rangle$B]}  
	 RIFIUTO: {'B32', 'A71', 'C01', 'B41', 'A77', 'A87', 'B33', 'C54'}\\ \tt PROPOSTA: {'B35', 'B32', 'A10', 'A71', 'C01', 'C54', 'B41', 'A77'}\\ \tt ARGOMENTAZIONE: {'Rimaniamo ampiamente sotto il limite: sforzo totale 4955 su 5138 (margine 183). Mantengo sei dei tuoi otto elementi (B32, A71, C01, B41, A77, C54) e propongo due sostituzioni mirate: A87 -> A10 e B33 -> B35. Il set risultante è più snello e coerente senza sacrificare il buffer di sforzo. Se per te è importante includere A87, posso scambiare B32 con A87 mantenendo invariato il resto (nuovo sforzo 5066). In alternativa, se preferisci introdurre C00, possiamo farlo 1:1 con B41 (nuovo sforzo 4944) oppure con C54 (nuovo sforzo 5085), lasciando inalterato tutto il resto. Dimmi quale di queste micro-varianti preferisci e chiudiamo rapidamente.'} 
	  } 
	   } 
	   } 
	 & \\ 
 

    \theutterance \stepcounter{utterance}  

    & & & \multicolumn{4}{p{0.6\linewidth}}{\cellcolor[rgb]{1,0.85,0.72}{%
	\makecell[{{p{\linewidth}}}]{% 
	  \tt {\tiny [GM$\langle$B]}  
	 RAGIONAMENTO STRATEGICO: {'La loro ultima proposta mi rende poco; il mio obiettivo minimo è far entrare A87 e B33 preservando i loro capisaldi dichiarati (A71, C01 e il pacchetto B41 + A77), ora già insieme a C54. La sostituzione mirata A10 -> A87 e B35 -> B33 porta il set a 4994 di sforzo (margine 144) e per me è un salto di valore netto importante. So che hanno appena rifiutato una versione simile, quindi offro anche due varianti di ripiego che rispettano le loro preferenze: (1) se vogliono tenere B35, posso scambiare B41 -> A04 mantenendo A87 e B33 (sforzo 4710, grande buffer); (2) se tengono ad A10, posso mantenere A10 e inserire comunque A87 e B33 rimuovendo B41 e B35 (sforzo 4967). Tenterò prima la mia combinazione principale perché allinea efficienza e stabilità e lascia ampio margine.'}\\ \tt \\ \tt RIFIUTO: {'B35', 'B32', 'A10', 'A71', 'C01', 'C54', 'B41', 'A77'}\\ \tt \\ \tt PROPOSTA: {'B32', 'A71', 'C01', 'C54', 'B41', 'A77', 'A87', 'B33'}\\ \tt \\ \tt ARGOMENTAZIONE: {'Resto ampiamente entro il limite: sforzo totale 4994 su 5138 (margine 144). Mantengo sei elementi della tua proposta (B32, A71, C01, C54, B41, A77) e applico due sostituzioni 1:1 ad alta efficienza: A10 -> A87 e B35 -> B33. Il mix risultante è più performante senza sacrificare il buffer. Se preferisci alternative: - Variante A (più buffer): teniamo B35 ma scambiamo B41 -> A04, mantenendo A87 e B33 (nuovo sforzo 4710). - Variante B (con A10): manteniamo A10 e inseriamo comunque A87 e B33 rimuovendo B41 e B35 (nuovo sforzo 4967). Dimmi quale configurazione senti più allineata e chiudiamo rapidamente.'} 
	  } 
	   } 
	   } 
	  \\ 
 

    \theutterance \stepcounter{utterance}  

    & & \multicolumn{4}{p{0.6\linewidth}}{\cellcolor[rgb]{0.9,0.9,0.9}{%
	\makecell[{{p{\linewidth}}}]{% 
	  \tt {\tiny [A$\langle$GM]}  
	 RIFIUTO: {'B35', 'B32', 'A10', 'A71', 'C01', 'C54', 'B41', 'A77'}\\ \tt \\ \tt PROPOSTA: {'B32', 'A71', 'C01', 'C54', 'B41', 'A77', 'A87', 'B33'}\\ \tt \\ \tt ARGOMENTAZIONE: {'Resto ampiamente entro il limite: sforzo totale 4994 su 5138 (margine 144). Mantengo sei elementi della tua proposta (B32, A71, C01, C54, B41, A77) e applico due sostituzioni 1:1 ad alta efficienza: A10 -> A87 e B35 -> B33. Il mix risultante è più performante senza sacrificare il buffer. Se preferisci alternative: - Variante A (più buffer): teniamo B35 ma scambiamo B41 -> A04, mantenendo A87 e B33 (nuovo sforzo 4710). - Variante B (con A10): manteniamo A10 e inseriamo comunque A87 e B33 rimuovendo B41 e B35 (nuovo sforzo 4967). Dimmi quale configurazione senti più allineata e chiudiamo rapidamente.'} 
	  } 
	   } 
	   } 
	 & \\ 
 

    \theutterance \stepcounter{utterance}  

    & \multicolumn{4}{p{0.6\linewidth}}{\cellcolor[rgb]{0.8,1,0.9}{%
	\makecell[{{p{\linewidth}}}]{% 
	  \tt {\tiny [A$\rangle$GM]}  
	 RAGIONAMENTO STRATEGICO: {'La loro proposta attuale mi rende poco (4138). Per favorire l’accordo con la minima modifica possibile, punto a un singolo scambio che alza molto il mio valore senza toccare i loro capisaldi: sostituire B33 con C00. Mantengo A87, A77, B41, e il nucleo B32-A71-C01-C54. Questa variante resta ben sotto soglia (4977) ed è la più redditizia tra i “one-swap”. Se B33 fosse irrinunciabile, piano B: tenere B33 e sostituire C54 con C00 (5124). Piano C: adottare la loro “Variante B” (B41 -> A10) a parità del resto (4967). Provo prima con lo swap B33->C00; se rifiutato, offrirò uno dei due piani di ripiego.'}\\ \tt RIFIUTO: {'B32', 'A71', 'C01', 'C54', 'B41', 'A77', 'A87', 'B33'}\\ \tt PROPOSTA: {'B32', 'A71', 'C01', 'C54', 'B41', 'A77', 'A87', 'C00'}\\ \tt ARGOMENTAZIONE: {'Restiamo ampiamente entro il limite: sforzo totale 4977 su 5138 (margine 161). Mantengo 7/8 degli elementi che proponi (B32, A71, C01, C54, B41, A77, A87) e applico un solo scambio mirato: B33 -> C00. Così il set resta snello e coerente con il mix che preferisci, ma guadagna solidità senza appesantire il budget. Se per te B33 è imprescindibile, posso tenere B33 e sostituire invece C54 con C00 (nuovo sforzo 5124) a parità del resto; in alternativa, posso adottare la tua “Variante B”: mantenere B33 e A87 ma scambiare B41 con A10 (nuovo sforzo 4967). Dimmi quale delle due micro-varianti senti più allineata: sono pronto a chiudere rapidamente.'} 
	  } 
	   } 
	   } 
	 & & \\ 
 

    \theutterance \stepcounter{utterance}  

    & & \multicolumn{4}{p{0.6\linewidth}}{\cellcolor[rgb]{0.9,0.9,0.9}{%
	\makecell[{{p{\linewidth}}}]{% 
	  \tt {\tiny [GM$\rangle$B]}  
	 RIFIUTO: {'B32', 'A71', 'C01', 'C54', 'B41', 'A77', 'A87', 'B33'}\\ \tt PROPOSTA: {'B32', 'A71', 'C01', 'C54', 'B41', 'A77', 'A87', 'C00'}\\ \tt ARGOMENTAZIONE: {'Restiamo ampiamente entro il limite: sforzo totale 4977 su 5138 (margine 161). Mantengo 7/8 degli elementi che proponi (B32, A71, C01, C54, B41, A77, A87) e applico un solo scambio mirato: B33 -> C00. Così il set resta snello e coerente con il mix che preferisci, ma guadagna solidità senza appesantire il budget. Se per te B33 è imprescindibile, posso tenere B33 e sostituire invece C54 con C00 (nuovo sforzo 5124) a parità del resto; in alternativa, posso adottare la tua “Variante B”: mantenere B33 e A87 ma scambiare B41 con A10 (nuovo sforzo 4967). Dimmi quale delle due micro-varianti senti più allineata: sono pronto a chiudere rapidamente.'} 
	  } 
	   } 
	   } 
	 & \\ 
 

    \theutterance \stepcounter{utterance}  

    & & & \multicolumn{4}{p{0.6\linewidth}}{\cellcolor[rgb]{1,0.85,0.72}{%
	\makecell[{{p{\linewidth}}}]{% 
	  \tt {\tiny [GM$\langle$B]}  
	 RAGIONAMENTO STRATEGICO: {'La loro proposta attuale mi dà valore 3686. Hanno però offerto esplicitamente la “Variante B” che include B33 e A10 al posto di B41, che per me è nettamente migliore (valore 3999) e resta comoda nel limite (4967). Per massimizzare le chance di accordo rapido, rilancio esattamente quella variante da loro suggerita. Come piano B, se preferiscono tenere C00, posso scambiare A10 -> C00 mantenendo invariato il resto (sforzo 4983), che per me resta accettabile.'}\\ \tt \\ \tt RIFIUTO: {'B32', 'A71', 'C01', 'C54', 'B41', 'A77', 'A87', 'C00'}\\ \tt \\ \tt PROPOSTA: {'B32', 'A71', 'C01', 'C54', 'A77', 'A87', 'B33', 'A10'}\\ \tt \\ \tt ARGOMENTAZIONE: {'Accolgo la “Variante B” che avevi proposto: restiamo entro il limite con sforzo totale 4967 su 5138 (margine 171). Manteniamo i tuoi capisaldi (A71, C01, A77) insieme a B32, C54 e A87, e sostituiamo B41 con A10 reinserendo B33 per un mix più efficiente a parità di budget. Se preferisci mantenere C00, posso scambiare A10 -> C00 a set invariato (nuovo sforzo 4983). Pronto a chiudere su questa configurazione.'} 
	  } 
	   } 
	   } 
	  \\ 
 

    \theutterance \stepcounter{utterance}  

    & & \multicolumn{4}{p{0.6\linewidth}}{\cellcolor[rgb]{0.9,0.9,0.9}{%
	\makecell[{{p{\linewidth}}}]{% 
	  \tt {\tiny [A$\langle$GM]}  
	 RIFIUTO: {'B32', 'A71', 'C01', 'C54', 'B41', 'A77', 'A87', 'C00'}\\ \tt \\ \tt PROPOSTA: {'B32', 'A71', 'C01', 'C54', 'A77', 'A87', 'B33', 'A10'}\\ \tt \\ \tt ARGOMENTAZIONE: {'Accolgo la “Variante B” che avevi proposto: restiamo entro il limite con sforzo totale 4967 su 5138 (margine 171). Manteniamo i tuoi capisaldi (A71, C01, A77) insieme a B32, C54 e A87, e sostituiamo B41 con A10 reinserendo B33 per un mix più efficiente a parità di budget. Se preferisci mantenere C00, posso scambiare A10 -> C00 a set invariato (nuovo sforzo 4983). Pronto a chiudere su questa configurazione.'} 
	  } 
	   } 
	   } 
	 & \\ 
 

    \theutterance \stepcounter{utterance}  

    & \multicolumn{4}{p{0.6\linewidth}}{\cellcolor[rgb]{0.8,1,0.9}{%
	\makecell[{{p{\linewidth}}}]{% 
	  \tt {\tiny [A$\rangle$GM]}  
	 RAGIONAMENTO STRATEGICO: {'La loro proposta attuale mi vale 4363. Hanno però già offerto esplicitamente lo scambio A10 -> C00, che per me migliora leggermente (4377) restando comodissimi nel limite (4983). Dato che hanno segnalato disponibilità a chiudere rapidamente su questa variante, punto a massimizzare la probabilità di accordo con la modifica minima già accettata da loro. Avrei preferito anche sostituire B33 con B35 (grande incremento per me), ma visto che lo hanno rifiutato più volte, evito di riaprire quel fronte e cerco di chiudere ora.'}\\ \tt RIFIUTO: {'B32', 'A71', 'C01', 'C54', 'A77', 'A87', 'B33', 'A10'}\\ \tt PROPOSTA: {'B32', 'A71', 'C01', 'C54', 'A77', 'A87', 'B33', 'C00'}\\ \tt ARGOMENTAZIONE: {'Accolgo la tua apertura: è la tua stessa configurazione con il solo scambio A10 -> C00 che avevi proposto come alternativa. Restiamo ampiamente entro il limite (sforzo totale 4983 su 5138, margine 155), manteniamo tutti gli altri elementi che desideri (A71, C01, C54, A77, A87, B33, B32) e chiudiamo senza ulteriori complicazioni. Se per qualsiasi motivo preferissi tornare ad A10, dimmelo: posso essere flessibile, ma questa variante è pronta per l’accordo immediato.'} 
	  } 
	   } 
	   } 
	 & & \\ 
 

    \theutterance \stepcounter{utterance}  

    & & \multicolumn{4}{p{0.6\linewidth}}{\cellcolor[rgb]{0.9,0.9,0.9}{%
	\makecell[{{p{\linewidth}}}]{% 
	  \tt {\tiny [GM$\rangle$B]}  
	 RIFIUTO: {'B32', 'A71', 'C01', 'C54', 'A77', 'A87', 'B33', 'A10'}\\ \tt PROPOSTA: {'B32', 'A71', 'C01', 'C54', 'A77', 'A87', 'B33', 'C00'}\\ \tt ARGOMENTAZIONE: {'Accolgo la tua apertura: è la tua stessa configurazione con il solo scambio A10 -> C00 che avevi proposto come alternativa. Restiamo ampiamente entro il limite (sforzo totale 4983 su 5138, margine 155), manteniamo tutti gli altri elementi che desideri (A71, C01, C54, A77, A87, B33, B32) e chiudiamo senza ulteriori complicazioni. Se per qualsiasi motivo preferissi tornare ad A10, dimmelo: posso essere flessibile, ma questa variante è pronta per l’accordo immediato.'} 
	  } 
	   } 
	   } 
	 & \\ 
 

    \theutterance \stepcounter{utterance}  

    & & & \multicolumn{4}{p{0.6\linewidth}}{\cellcolor[rgb]{1,0.85,0.72}{%
	\makecell[{{p{\linewidth}}}]{% 
	  \tt {\tiny [GM$\langle$B]}  
	 RAGIONAMENTO STRATEGICO: {'La loro proposta corrente mi dà 3982 di valore. Con un’unica sostituzione C00 -> A04, il mio valore sale di molto (+701) mantenendo invariati 7/8 elementi che a loro piacciono, e riduco lo sforzo di 273 (4983 -> 4710), aumentando il buffer: è una modifica minima, coerente con la loro enfasi su set “snelli” e margine di sicurezza. Se spingono per tenere C00, posso comunque chiudere sul loro set attuale; in alternativa, potrei accettare una variante che mantenga C00 inserendo A04 al posto di A71 (so che per loro A71 ricorre spesso, quindi la propongo solo come opzione B). Tenterò prima lo scambio secco C00 -> A04 per massimizzare il mio valore con il minimo cambiamento.'}\\ \tt \\ \tt RIFIUTO: {'B32', 'A71', 'C01', 'C54', 'A77', 'A87', 'B33', 'C00'}\\ \tt \\ \tt PROPOSTA: {'B32', 'A71', 'C01', 'C54', 'A77', 'A87', 'B33', 'A04'}\\ \tt \\ \tt ARGOMENTAZIONE: {'Restiamo comodamente entro il limite: sforzo totale 4710 su 5138 (margine 428). Tengo 7 degli 8 elementi della tua proposta e applico un solo scambio mirato C00 -> A04: così il set diventa ancora più snello (−273 di sforzo) e con maggiore impatto, senza toccare gli altri componenti che hai richiesto. Se per te è imprescindibile mantenere C00, posso chiudere sulla tua proposta attuale; in alternativa, possiamo mantenere C00 e inserire A04 sostituendo A71 (sforzo 4606), lasciando invariato il resto. Dimmi quale preferisci e chiudiamo subito.'} 
	  } 
	   } 
	   } 
	  \\ 
 

    \theutterance \stepcounter{utterance}  

    & & \multicolumn{4}{p{0.6\linewidth}}{\cellcolor[rgb]{0.9,0.9,0.9}{%
	\makecell[{{p{\linewidth}}}]{% 
	  \tt {\tiny [A$\langle$GM]}  
	 RIFIUTO: {'B32', 'A71', 'C01', 'C54', 'A77', 'A87', 'B33', 'C00'}\\ \tt \\ \tt PROPOSTA: {'B32', 'A71', 'C01', 'C54', 'A77', 'A87', 'B33', 'A04'}\\ \tt \\ \tt ARGOMENTAZIONE: {'Restiamo comodamente entro il limite: sforzo totale 4710 su 5138 (margine 428). Tengo 7 degli 8 elementi della tua proposta e applico un solo scambio mirato C00 -> A04: così il set diventa ancora più snello (−273 di sforzo) e con maggiore impatto, senza toccare gli altri componenti che hai richiesto. Se per te è imprescindibile mantenere C00, posso chiudere sulla tua proposta attuale; in alternativa, possiamo mantenere C00 e inserire A04 sostituendo A71 (sforzo 4606), lasciando invariato il resto. Dimmi quale preferisci e chiudiamo subito.'} 
	  } 
	   } 
	   } 
	 & \\ 
 

    \theutterance \stepcounter{utterance}  

    & \multicolumn{4}{p{0.6\linewidth}}{\cellcolor[rgb]{0.8,1,0.9}{%
	\makecell[{{p{\linewidth}}}]{% 
	  \tt {\tiny [A$\rangle$GM]}  
	 RAGIONAMENTO STRATEGICO: {'La loro nuova proposta con A04 riduce molto il mio rendimento, mentre hanno appena dichiarato disponibilità a chiudere sulla variante con C00 che avevo proposto. Per massimizzare le chance di accordo senza riaprire molte modifiche, rifiuto A04 e ripropongo esattamente il set con C00, che resta comodo nel budget ed è quello su cui hanno già offerto chiusura. Come piano B, se volessero evitare C00, potrei accettare la stessa configurazione con A10 al posto di C00, ma provo prima a consolidare C00.'}\\ \tt RIFIUTO: {'B32', 'A71', 'C01', 'C54', 'A77', 'A87', 'B33', 'A04'}\\ \tt PROPOSTA: {'B32', 'A71', 'C01', 'C54', 'A77', 'A87', 'B33', 'C00'}\\ \tt ARGOMENTAZIONE: {'Restiamo ampiamente entro il limite: sforzo totale 4983 su 5138 (margine 155). Mantengo 7/8 elementi della tua proposta e applico un solo scambio mirato A04 -> C00, come avevi già indicato di poter accettare. Il set è snello, coerente con quanto hai richiesto e pronto per chiudere subito. Se preferisci evitare C00, posso passare alla stessa configurazione con A10 al posto di C00 mantenendo invariato tutto il resto; dimmi solo quale delle due varianti confermi.'} 
	  } 
	   } 
	   } 
	 & & \\ 
 

    \theutterance \stepcounter{utterance}  

    & & \multicolumn{4}{p{0.6\linewidth}}{\cellcolor[rgb]{0.9,0.9,0.9}{%
	\makecell[{{p{\linewidth}}}]{% 
	  \tt {\tiny [GM$\rangle$B]}  
	 RIFIUTO: {'B32', 'A71', 'C01', 'C54', 'A77', 'A87', 'B33', 'A04'}\\ \tt PROPOSTA: {'B32', 'A71', 'C01', 'C54', 'A77', 'A87', 'B33', 'C00'}\\ \tt ARGOMENTAZIONE: {'Restiamo ampiamente entro il limite: sforzo totale 4983 su 5138 (margine 155). Mantengo 7/8 elementi della tua proposta e applico un solo scambio mirato A04 -> C00, come avevi già indicato di poter accettare. Il set è snello, coerente con quanto hai richiesto e pronto per chiudere subito. Se preferisci evitare C00, posso passare alla stessa configurazione con A10 al posto di C00 mantenendo invariato tutto il resto; dimmi solo quale delle due varianti confermi.'} 
	  } 
	   } 
	   } 
	 & \\ 
 

    \theutterance \stepcounter{utterance}  

    & & & \multicolumn{4}{p{0.6\linewidth}}{\cellcolor[rgb]{1,0.85,0.72}{%
	\makecell[{{p{\linewidth}}}]{% 
	  \tt {\tiny [GM$\langle$B]}  
	 RAGIONAMENTO STRATEGICO: {'La tua variante con A10 al posto di C00 mi offre un valore leggermente superiore (e sforzo più basso) rispetto a C00, restando identica per 7/8 del set che preferisci. Tento quindi la chiusura sulla tua stessa alternativa A10; se per te è essenziale C00, sono disposto ad accettare immediatamente anche la tua ultima proposta con C00 pur di chiudere.'}\\ \tt \\ \tt PROPOSTA: {'B32', 'A71', 'C01', 'C54', 'A77', 'A87', 'B33', 'A10'}\\ \tt \\ \tt ARGOMENTAZIONE: {'Questa è esattamente la tua alternativa con A10 al posto di C00: sforzo totale 4967 su 5138 (margine 171). Manteniamo invariati tutti gli altri elementi che hai richiesto (B32, A71, C01, C54, A77, A87, B33). Se preferisci tenere C00, posso confermare subito anche la tua ultima proposta senza ulteriori modifiche; dimmi quale delle due chiudiamo.'} 
	  } 
	   } 
	   } 
	  \\ 
 

    \theutterance \stepcounter{utterance}  

    & & \multicolumn{4}{p{0.6\linewidth}}{\cellcolor[rgb]{0.9,0.9,0.9}{%
	\makecell[{{p{\linewidth}}}]{% 
	  \tt {\tiny [A$\langle$GM]}  
	 PROPOSTA: {'B32', 'A71', 'C01', 'C54', 'A77', 'A87', 'B33', 'A10'}\\ \tt \\ \tt ARGOMENTAZIONE: {'Questa è esattamente la tua alternativa con A10 al posto di C00: sforzo totale 4967 su 5138 (margine 171). Manteniamo invariati tutti gli altri elementi che hai richiesto (B32, A71, C01, C54, A77, A87, B33). Se preferisci tenere C00, posso confermare subito anche la tua ultima proposta senza ulteriori modifiche; dimmi quale delle due chiudiamo.'} 
	  } 
	   } 
	   } 
	 & \\ 
 

    \theutterance \stepcounter{utterance}  

    & \multicolumn{4}{p{0.6\linewidth}}{\cellcolor[rgb]{0.8,1,0.9}{%
	\makecell[{{p{\linewidth}}}]{% 
	  \tt {\tiny [A$\rangle$GM]}  
	 RAGIONAMENTO STRATEGICO: {'La differenza per me tra A10 e C00 è minima (+14 per C00), mentre spingere ancora rischia di far saltare l’accordo. Hanno già rifiutato più volte le mie richieste chiave (B35 al posto di B33, reintroduzione stabile di C00) e ora offrono una chiusura pulita su una variante che è comunque solida per me. Accetto per massimizzare la probabilità di accordo immediato con valore stabile.'}\\ \tt ACCORDO: {'B32', 'A71', 'C01', 'C54', 'A77', 'A87', 'B33', 'A10'}\\ \tt ARGOMENTAZIONE: {'Accetto la tua proposta: restiamo entro il limite (sforzo totale 4967 su 5138) e manteniamo tutti gli elementi che hai indicato. Grazie per la collaborazione: chiudiamo su questo set.'} 
	  } 
	   } 
	   } 
	 & & \\ 
 

    \theutterance \stepcounter{utterance}  

    & & & \multicolumn{2}{p{0.3\linewidth}}{\cellcolor[rgb]{0.95,0.95,0.95}{%
	\makecell[{{p{\linewidth}}}]{% 
	  \tt {\tiny [GM$|$GM]}  
	 game successful 
	  } 
	   } 
	   } 
	 & & \\ 
 

    \theutterance \stepcounter{utterance}  

    & & & \multicolumn{2}{p{0.3\linewidth}}{\cellcolor[rgb]{0.95,0.95,0.95}{%
	\makecell[{{p{\linewidth}}}]{% 
	  \tt {\tiny [GM$|$GM]}  
	 end game 
	  } 
	   } 
	   } 
	 & & \\ 
 

\end{supertabular}
}

\end{document}
